\begingroup
\small
\setlength{\parskip}{.12\baselineskip}
\setlist[itemize]{leftmargin=1.65em,itemsep=0pt,topsep=.1em,
  parsep=0pt,partopsep=0pt}
\let\texttt\textnormal
Las fichas siguientes publican el operador en su codominio nativo. Cada espectro se reproduce íntegramente; ninguna lista de autovalores se sustituye por una coordenada de masa.
\Needspace{10\baselineskip}
\noindent\textbf{sector leptónico cargado, profundidad G0 (centro).} \(\mathcal B_{\ell,\mathrm{G0}}\).\par
\noindent\textbf{Dominio de fibra.} Rango 1; celda 41; sin órbita cromática; 1 familia de ruta.
\noindent\textbf{Espectro completo de \(K_D\).}
\begin{itemize}
\item \texttt{\detokenize{(80,54,6)x1}}
\end{itemize}
\noindent\textbf{Espectro completo de \(K_\Omega\).}
\begin{itemize}
\item \texttt{\detokenize{(80,54,6)x1}}
\end{itemize}
\noindent\textbf{Conclusión en la fibra.} Operador diagonal exacto sobre la fibra \(B_1\); la ruta nominal y el morfismo declarado publican el carácter escalar sin sustituir el espectro de la fibra.
\medskip
\Needspace{10\baselineskip}
\noindent\textbf{sector leptónico cargado, profundidad G2.} \(\mathcal B_{\ell,\mathrm{G2}}\).\par
\noindent\textbf{Dominio de fibra.} Rango 8; celdas 21,23,25,39,43,57,59,61; sin órbita cromática; 8 familias de ruta.
\noindent\textbf{Espectro completo de \(K_D\).}
\begin{itemize}
\item \texttt{\detokenize{(0,24,0)x1}}
\item \texttt{\detokenize{(40,78,27)x2}}
\item \texttt{\detokenize{(80,54,7)x1}}
\item \texttt{\detokenize{(80,66,2)x1}}
\item \texttt{\detokenize{(80,66,3)x1}}
\item \texttt{\detokenize{(100,66,0)x1}}
\item \texttt{\detokenize{(100,66,2)x1}}
\end{itemize}
\noindent\textbf{Espectro completo de \(K_\Omega\).}
\begin{itemize}
\item \texttt{\detokenize{(80,54,6)x8}}
\end{itemize}
\noindent\textbf{Conclusión en la fibra.} Operador diagonal exacto sobre la fibra \(B_1\); la ruta nominal y el morfismo declarado publican el carácter escalar sin sustituir el espectro de la fibra.
\medskip
\Needspace{22\baselineskip}
\noindent\textbf{sector leptónico cargado, profundidad G4.} \(\mathcal B_{\ell,\mathrm{G4}}\).\par
\noindent\textbf{Dominio de fibra.} Rango 16; celdas 1,3,5,7,9,19,27,37,45,55,63,73,75,77,79,81; sin órbita cromática; 16 familias de ruta.
\noindent\textbf{Espectro completo de \(K_D\).}
\begin{itemize}
\item \texttt{\detokenize{(0,24,0)x1}}
\item \texttt{\detokenize{(40,24,2)x2}}
\item \texttt{\detokenize{(40,54,6)x2}}
\item \texttt{\detokenize{(40,84,30)x2}}
\item \texttt{\detokenize{(60,78,25)x3}}
\item \texttt{\detokenize{(80,66,2)x2}}
\item \texttt{\detokenize{(80,66,3)x3}}
\item \texttt{\detokenize{(80,66,4)x1}}
\end{itemize}
\noindent\textbf{Espectro completo de \(K_\Omega\).}
\begin{itemize}
\item \texttt{\detokenize{(24,54,2)x1}}
\item \texttt{\detokenize{(56,54,5)x4}}
\item \texttt{\detokenize{(80,54,6)x11}}
\end{itemize}
\noindent\textbf{Conclusión en la fibra.} Operador diagonal exacto sobre la fibra \(B_1\); la ruta nominal y el morfismo declarado publican el carácter escalar sin sustituir el espectro de la fibra.
\medskip
\Needspace{10\baselineskip}
\noindent\textbf{sector neutrínico, profundidad G1.} \(\mathcal B_{\nu,\mathrm{G1}}\).\par
\noindent\textbf{Dominio de fibra.} Rango 2; celdas 40,42; sin órbita cromática; 2 familias de ruta.
\noindent\textbf{Espectro completo de \(K_D\).}
\begin{itemize}
\item \texttt{\detokenize{(40,24,2)x1}}
\item \texttt{\detokenize{(40,78,27)x1}}
\end{itemize}
\noindent\textbf{Espectro completo de \(K_\Omega\).}
\begin{itemize}
\item \texttt{\detokenize{(40,54,4)x1}}
\item \texttt{\detokenize{(80,54,6)x1}}
\end{itemize}
\noindent\textbf{Conclusión en la fibra.} Operador diagonal exacto sobre la fibra \(B_1\); la ruta nominal y el morfismo declarado publican el carácter escalar sin sustituir el espectro de la fibra.
\medskip
\Needspace{10\baselineskip}
\noindent\textbf{sector neutrínico, profundidad G2.} \(\mathcal B_{\nu,\mathrm{G2}}\).\par
\noindent\textbf{Dominio de fibra.} Rango 4; celdas 22,24,58,60; sin órbita cromática; 4 familias de ruta.
\noindent\textbf{Espectro completo de \(K_D\).}
\begin{itemize}
\item \texttt{\detokenize{(0,24,0)x1}}
\item \texttt{\detokenize{(40,84,30)x1}}
\item \texttt{\detokenize{(60,78,25)x1}}
\item \texttt{\detokenize{(80,66,3)x1}}
\end{itemize}
\noindent\textbf{Espectro completo de \(K_\Omega\).}
\begin{itemize}
\item \texttt{\detokenize{(4,54,2)x1}}
\item \texttt{\detokenize{(80,54,6)x3}}
\end{itemize}
\noindent\textbf{Conclusión en la fibra.} Operador diagonal exacto sobre la fibra \(B_1\); la ruta nominal y el morfismo declarado publican el carácter escalar sin sustituir el espectro de la fibra.
\medskip
\Needspace{18\baselineskip}
\noindent\textbf{sector neutrínico, profundidad G3.} \(\mathcal B_{\nu,\mathrm{G3}}\).\par
\noindent\textbf{Dominio de fibra.} Rango 6; celdas 20,26,38,44,56,62; sin órbita cromática; 6 familias de ruta.
\noindent\textbf{Espectro completo de \(K_D\).}
\begin{itemize}
\item \texttt{\detokenize{(40,24,2)x2}}
\item \texttt{\detokenize{(40,78,27)x1}}
\item \texttt{\detokenize{(60,84,28)x1}}
\item \texttt{\detokenize{(80,54,6)x1}}
\item \texttt{\detokenize{(80,66,3)x1}}
\end{itemize}
\noindent\textbf{Espectro completo de \(K_\Omega\).}
\begin{itemize}
\item \texttt{\detokenize{(36,54,4)x1}}
\item \texttt{\detokenize{(56,54,5)x1}}
\item \texttt{\detokenize{(80,54,6)x4}}
\end{itemize}
\noindent\textbf{Conclusión en la fibra.} Operador diagonal exacto sobre la fibra \(B_1\); la ruta nominal y el morfismo declarado publican el carácter escalar sin sustituir el espectro de la fibra.
\medskip
\Needspace{10\baselineskip}
\noindent\textbf{sector neutrínico, profundidad G4.} \(\mathcal B_{\nu,\mathrm{G4}}\).\par
\noindent\textbf{Dominio de fibra.} Rango 8; celdas 2,4,6,8,74,76,78,80; sin órbita cromática; 8 familias de ruta.
\noindent\textbf{Espectro completo de \(K_D\).}
\begin{itemize}
\item \texttt{\detokenize{(0,24,0)x1}}
\item \texttt{\detokenize{(40,24,2)x1}}
\item \texttt{\detokenize{(40,78,27)x2}}
\item \texttt{\detokenize{(40,84,30)x1}}
\item \texttt{\detokenize{(80,54,7)x1}}
\item \texttt{\detokenize{(80,66,2)x1}}
\item \texttt{\detokenize{(100,66,0)x1}}
\end{itemize}
\noindent\textbf{Espectro completo de \(K_\Omega\).}
\begin{itemize}
\item \texttt{\detokenize{(36,54,4)x1}}
\item \texttt{\detokenize{(56,54,5)x1}}
\item \texttt{\detokenize{(80,54,6)x6}}
\end{itemize}
\noindent\textbf{Conclusión en la fibra.} Operador diagonal exacto sobre la fibra \(B_1\); la ruta nominal y el morfismo declarado publican el carácter escalar sin sustituir el espectro de la fibra.
\medskip
\Needspace{10\baselineskip}
\noindent\textbf{rama quark de tipo inferior, profundidad G1.} \(\mathcal B_{d,\mathrm{G1}}\).\par
\noindent\textbf{Dominio de fibra.} Rango 2; celdas 32,50; B | G; 2 familias de ruta.
\noindent\textbf{Espectro completo de \(K_D\).}
\begin{itemize}
\item \texttt{\detokenize{(40,24,2)x1}}
\item \texttt{\detokenize{(40,78,27)x1}}
\end{itemize}
\noindent\textbf{Espectro completo de \(K_\Omega\).}
\begin{itemize}
\item \texttt{\detokenize{(40,54,4)x1}}
\item \texttt{\detokenize{(80,54,6)x1}}
\end{itemize}
\noindent\textbf{Conclusión en la fibra.} Operador diagonal exacto sobre la fibra \(B_1\); la ruta nominal y el morfismo declarado publican el carácter escalar sin sustituir el espectro de la fibra.
\medskip
\Needspace{10\baselineskip}
\noindent\textbf{rama quark de tipo inferior, profundidad G2.} \(\mathcal B_{d,\mathrm{G2}}\).\par
\noindent\textbf{Dominio de fibra.} Rango 4; celdas 30,34,48,52; B | G | R; 4 familias de ruta.
\noindent\textbf{Espectro completo de \(K_D\).}
\begin{itemize}
\item \texttt{\detokenize{(0,24,0)x1}}
\item \texttt{\detokenize{(40,84,30)x1}}
\item \texttt{\detokenize{(60,78,25)x1}}
\item \texttt{\detokenize{(80,66,2)x1}}
\end{itemize}
\noindent\textbf{Espectro completo de \(K_\Omega\).}
\begin{itemize}
\item \texttt{\detokenize{(4,54,2)x1}}
\item \texttt{\detokenize{(80,54,6)x3}}
\end{itemize}
\noindent\textbf{Conclusión en la fibra.} Operador diagonal exacto sobre la fibra \(B_1\); la ruta nominal y el morfismo declarado publican el carácter escalar sin sustituir el espectro de la fibra.
\medskip
\Needspace{10\baselineskip}
\noindent\textbf{rama quark de tipo inferior, profundidad G3.} \(\mathcal B_{d,\mathrm{G3}}\).\par
\noindent\textbf{Dominio de fibra.} Rango 6; celdas 12,14,16,66,68,70; B | G | R; 6 familias de ruta.
\noindent\textbf{Espectro completo de \(K_D\).}
\begin{itemize}
\item \texttt{\detokenize{(40,24,2)x2}}
\item \texttt{\detokenize{(40,78,27)x1}}
\item \texttt{\detokenize{(60,84,28)x1}}
\item \texttt{\detokenize{(80,54,6)x1}}
\item \texttt{\detokenize{(80,66,4)x1}}
\end{itemize}
\noindent\textbf{Espectro completo de \(K_\Omega\).}
\begin{itemize}
\item \texttt{\detokenize{(36,54,4)x1}}
\item \texttt{\detokenize{(80,54,6)x5}}
\end{itemize}
\noindent\textbf{Conclusión en la fibra.} Operador diagonal exacto sobre la fibra \(B_1\); la ruta nominal y el morfismo declarado publican el carácter escalar sin sustituir el espectro de la fibra.
\medskip
\Needspace{10\baselineskip}
\noindent\textbf{rama quark de tipo inferior, profundidad G4.} \(\mathcal B_{d,\mathrm{G4}}\).\par
\noindent\textbf{Dominio de fibra.} Rango 8; celdas 10,18,28,36,46,54,64,72; B | G | R; 8 familias de ruta.
\noindent\textbf{Espectro completo de \(K_D\).}
\begin{itemize}
\item \texttt{\detokenize{(0,24,0)x1}}
\item \texttt{\detokenize{(40,24,2)x1}}
\item \texttt{\detokenize{(40,78,27)x2}}
\item \texttt{\detokenize{(40,84,30)x1}}
\item \texttt{\detokenize{(80,54,7)x1}}
\item \texttt{\detokenize{(80,66,3)x1}}
\item \texttt{\detokenize{(100,66,2)x1}}
\end{itemize}
\noindent\textbf{Espectro completo de \(K_\Omega\).}
\begin{itemize}
\item \texttt{\detokenize{(36,54,4)x1}}
\item \texttt{\detokenize{(80,54,6)x7}}
\end{itemize}
\noindent\textbf{Conclusión en la fibra.} Operador diagonal exacto sobre la fibra \(B_1\); la ruta nominal y el morfismo declarado publican el carácter escalar sin sustituir el espectro de la fibra.
\medskip
\Needspace{10\baselineskip}
\noindent\textbf{rama quark de tipo superior, profundidad G1.} \(\mathcal B_{u,\mathrm{G1}}\).\par
\noindent\textbf{Dominio de fibra.} Rango 4; celdas 31,33,49,51; B | G | R; 3 familias de ruta.
\noindent\textbf{Espectro completo de \(K_D\).}
\begin{itemize}
\item \texttt{\detokenize{(40,54,6)x1}}
\item \texttt{\detokenize{(60,78,25)x1}}
\item \texttt{\detokenize{(80,66,2)x1}}
\item \texttt{\detokenize{(80,66,3)x1}}
\end{itemize}
\noindent\textbf{Espectro completo de \(K_\Omega\).}
\begin{itemize}
\item \texttt{\detokenize{(80,54,6)x4}}
\end{itemize}
\noindent\textbf{Conclusión en la fibra.} Operador diagonal exacto sobre la fibra \(B_1\); la ruta nominal y el morfismo declarado publican el carácter escalar sin sustituir el espectro de la fibra.
\medskip
\Needspace{29\baselineskip}
\noindent\textbf{rama quark de tipo superior, profundidad G3.} \(\mathcal B_{u,\mathrm{G3}}\).\par
\noindent\textbf{Dominio de fibra.} Rango 12; celdas 11,13,15,17,29,35,47,53,65,67,69,71; B | G | R; 12 familias de ruta.
\noindent\textbf{Espectro completo de \(K_D\).}
\begin{itemize}
\item \texttt{\detokenize{(40,24,2)x2}}
\item \texttt{\detokenize{(40,78,27)x2}}
\item \texttt{\detokenize{(60,84,28)x1}}
\item \texttt{\detokenize{(80,54,6)x3}}
\item \texttt{\detokenize{(80,66,3)x1}}
\item \texttt{\detokenize{(80,66,4)x1}}
\item \texttt{\detokenize{(100,66,0)x1}}
\item \texttt{\detokenize{(100,66,2)x1}}
\end{itemize}
\noindent\textbf{Espectro completo de \(K_\Omega\).}
\begin{itemize}
\item \texttt{\detokenize{(56,54,5)x3}}
\item \texttt{\detokenize{(80,54,6)x9}}
\end{itemize}
\noindent\textbf{Conclusión en la fibra.} Operador diagonal exacto sobre la fibra \(B_1\); la ruta nominal y el morfismo declarado publican el carácter escalar sin sustituir el espectro de la fibra.
\medskip
\endgroup
